
\documentclass{article}
\usepackage{colortbl}
\usepackage{makecell}
\usepackage{multirow}
\usepackage{supertabular}

\begin{document}

\newcounter{utterance}

\centering \large Interaction Transcript for game `dond', experiment `semi\_en', episode 13 with Qwen3{-}32B{-}t0.0.
\vspace{24pt}

{ \footnotesize  \setcounter{utterance}{1}
\setlength{\tabcolsep}{0pt}
\begin{supertabular}{c@{$\;$}|p{.15\linewidth}@{}p{.15\linewidth}p{.15\linewidth}p{.15\linewidth}p{.15\linewidth}p{.15\linewidth}}
    \# & $\;$A & \multicolumn{4}{c}{Game Master} & $\;\:$B\\
    \hline

    \theutterance \stepcounter{utterance}  
    & & \multicolumn{4}{p{0.6\linewidth}}{
        \cellcolor[rgb]{0.9,0.9,0.9}{
            \makecell[{{p{\linewidth}}}]{
                \texttt{\tiny{[P1$\langle$GM]}}
                \texttt{You are playing a negotiation game in which you have to agree on how to divide a set of items among you and another player.} \\
\\ 
\texttt{Rules:} \\
\texttt{(a) You and the other player are given a set of items. Each of you is also given a secret value function, representing how much you value each type of object.} \\
\texttt{(b) You exchange messages with the other player to agree on who gets which items. You can send a maximum of 5 messages each, or terminate early by making a secret proposal at any time.} \\
\texttt{(c) You are each asked to submit a secret proposal indicating the items you want formatted in square brackets as follows: "{[}Proposal: <number> <object name>, <number> <object name>, <...>{]}"} \\
\texttt{(d) If your proposals are complementary, i.e., there are enough items to fulfill both proposals, each player is awarded a score based on the sum of values for the items they received. Otherwise, both of you get zero points.} \\
\texttt{(e) Your goal is to maximize the score you receive.} \\
\\ 
\texttt{Let us start.} \\
\\ 
\texttt{The set of available items is:} \\
\\ 
\texttt{4 papers, 2 clocks, 2 potatoes} \\
\\ 
\texttt{Your secret value function is:} \\
\\ 
\texttt{paper: 1, clock: 1, potato: 2} \\
\\ 
\texttt{IMPORTANT: Your messages, unless it is the secret proposal, are directly transmitted to the other player, so do not include any response to the rules or text announcing your message. To make a secret proposal, use the indicated format. Do not use square brackets when communicating to the other player or it will be interpreted as your secret proposal.} \\
            }
        }
    }
    & \\ \\

    \theutterance \stepcounter{utterance}  
    & \multicolumn{4}{p{0.6\linewidth}}{
        \cellcolor[rgb]{0.9,0.9,0.9}{
            \makecell[{{p{\linewidth}}}]{
                \texttt{\tiny{[P1$\rangle$GM]}}
                \texttt{<think>} \\
\texttt{Okay, let's see. I need to negotiate the division of 4 papers, 2 clocks, and 2 potatoes with another player. My value function is paper:1, clock:1, potato:2. So, I value potatoes the most, then papers and clocks equally. The other player's value function is secret, so I don't know what they want. The goal is to maximize my score, which is the sum of the values of the items I get. But if our proposals aren't complementary, we both get zero. So I need to figure out a way to propose a division that both of us can agree on.} \\
\\ 
\texttt{First, I should consider what the other player might value. Since I don't know their values, maybe I can make an offer that's fair or that leaves enough for them. Let's think about possible scenarios. If the other player also values potatoes highly, maybe they want some. But if they value papers or clocks more, they might prefer those. Since there are 2 potatoes, maybe I can take both, but then the other player might need some. Alternatively, I could take one potato and offer the other to them. But how to structure the proposal?} \\
\\ 
\texttt{My value function gives me 2 per potato, so each potato is worth more to me than a paper or clock. So ideally, I want as many potatoes as possible. But there are only two. Let's see. If I take both potatoes, that's 4 points. Then the remaining items are 4 papers and 2 clocks. The other player might want some of those. But if they need, say, 2 papers and 1 clock, that's possible. But I need to make sure that my proposal doesn't conflict with theirs. For example, if I take both potatoes, and they also want both potatoes, then our proposals would conflict, and we both get zero. So maybe I should take one potato and offer the other to them, but then I can get some papers or clocks as well.} \\
\\ 
\texttt{Alternatively, maybe I can propose taking 1 potato and 2 papers. That would give me 2 (from potato) + 2 (from 2 papers) = 4 points. Then the other player could take the remaining 1 potato, 2 papers, and 2 clocks. But I don't know their values. If they value clocks more, maybe they want the clocks. But if they take the clocks} \\
            }
        }
    }
    & & \\ \\

    \theutterance \stepcounter{utterance}  
    & & & \multicolumn{2}{p{0.3\linewidth}}{
        \cellcolor[rgb]{0.9,0.9,0.9}{
            \makecell[{{p{\linewidth}}}]{
                \texttt{\tiny{[GM$|$GM]}}
                \texttt{continue} \\
            }
        }
    }
    & & \\ \\

    \theutterance \stepcounter{utterance}  
    & & & \multicolumn{2}{p{0.3\linewidth}}{
        \cellcolor[rgb]{0.9,0.9,0.9}{
            \makecell[{{p{\linewidth}}}]{
                \texttt{\tiny{[GM$|$GM]}}
                \texttt{<think>} \\
\texttt{Okay, let's see. I need to negotiate the division of 4 papers, 2 clocks, and 2 potatoes with another player. My value function is paper:1, clock:1, potato:2. So, I value potatoes the most, then papers and clocks equally. The other player's value function is secret, so I don't know what they want. The goal is to maximize my score, which is the sum of the values of the items I get. But if our proposals aren't complementary, we both get zero. So I need to figure out a way to propose a division that both of us can agree on.} \\
\\ 
\texttt{First, I should consider what the other player might value. Since I don't know their values, maybe I can make an offer that's fair or that leaves enough for them. Let's think about possible scenarios. If the other player also values potatoes highly, maybe they want some. But if they value papers or clocks more, they might prefer those. Since there are 2 potatoes, maybe I can take both, but then the other player might need some. Alternatively, I could take one potato and offer the other to them. But how to structure the proposal?} \\
\\ 
\texttt{My value function gives me 2 per potato, so each potato is worth more to me than a paper or clock. So ideally, I want as many potatoes as possible. But there are only two. Let's see. If I take both potatoes, that's 4 points. Then the remaining items are 4 papers and 2 clocks. The other player might want some of those. But if they need, say, 2 papers and 1 clock, that's possible. But I need to make sure that my proposal doesn't conflict with theirs. For example, if I take both potatoes, and they also want both potatoes, then our proposals would conflict, and we both get zero. So maybe I should take one potato and offer the other to them, but then I can get some papers or clocks as well.} \\
\\ 
\texttt{Alternatively, maybe I can propose taking 1 potato and 2 papers. That would give me 2 (from potato) + 2 (from 2 papers) = 4 points. Then the other player could take the remaining 1 potato, 2 papers, and 2 clocks. But I don't know their values. If they value clocks more, maybe they want the clocks. But if they take the clocks} \\
            }
        }
    }
    & & \\ \\

    \theutterance \stepcounter{utterance}  
    & & \multicolumn{4}{p{0.6\linewidth}}{
        \cellcolor[rgb]{0.9,0.9,0.9}{
            \makecell[{{p{\linewidth}}}]{
                \texttt{\tiny{[GM$\rangle$P2]}}
                \texttt{You are playing a negotiation game in which you have to agree on how to divide a set of items among you and another player.} \\
\\ 
\texttt{Rules:} \\
\texttt{(a) You and the other player are given a set of items. Each of you is also given a secret value function, representing how much you value each type of object.} \\
\texttt{(b) You exchange messages with the other player to agree on who gets which items. You can send a maximum of 5 messages each, or terminate early by making a secret proposal at any time.} \\
\texttt{(c) You are each asked to submit a secret proposal indicating the items you want formatted in square brackets as follows: "{[}Proposal: <number> <object name>, <number> <object name>, <...>{]}"} \\
\texttt{(d) If your proposals are complementary, i.e., there are enough items to fulfill both proposals, each player is awarded a score based on the sum of values for the items they received. Otherwise, both of you get zero points.} \\
\texttt{(e) Your goal is to maximize the score you receive.} \\
\\ 
\texttt{Let us start.} \\
\\ 
\texttt{The set of available items is:} \\
\\ 
\texttt{4 papers, 2 clocks, 2 potatoes} \\
\\ 
\texttt{Your secret value function is:} \\
\\ 
\texttt{paper: 0, clock: 4, potato: 1} \\
\\ 
\texttt{IMPORTANT: Your messages, unless it is the secret proposal, are directly transmitted to the other player, so do not include any response to the rules or text announcing your message. To make a secret proposal, use the indicated format. Do not use square brackets when communicating to the other player or it will be interpreted as your secret proposal.} \\
\\ 
\\ 
\texttt{<think>} \\
\texttt{Okay, let's see. I need to negotiate the division of 4 papers, 2 clocks, and 2 potatoes with another player. My value function is paper:1, clock:1, potato:2. So, I value potatoes the most, then papers and clocks equally. The other player's value function is secret, so I don't know what they want. The goal is to maximize my score, which is the sum of the values of the items I get. But if our proposals aren't complementary, we both get zero. So I need to figure out a way to propose a division that both of us can agree on.} \\
\\ 
\texttt{First, I should consider what the other player might value. Since I don't know their values, maybe I can make an offer that's fair or that leaves enough for them. Let's think about possible scenarios. If the other player also values potatoes highly, maybe they want some. But if they value papers or clocks more, they might prefer those. Since there are 2 potatoes, maybe I can take both, but then the other player might need some. Alternatively, I could take one potato and offer the other to them. But how to structure the proposal?} \\
\\ 
\texttt{My value function gives me 2 per potato, so each potato is worth more to me than a paper or clock. So ideally, I want as many potatoes as possible. But there are only two. Let's see. If I take both potatoes, that's 4 points. Then the remaining items are 4 papers and 2 clocks. The other player might want some of those. But if they need, say, 2 papers and 1 clock, that's possible. But I need to make sure that my proposal doesn't conflict with theirs. For example, if I take both potatoes, and they also want both potatoes, then our proposals would conflict, and we both get zero. So maybe I should take one potato and offer the other to them, but then I can get some papers or clocks as well.} \\
\\ 
\texttt{Alternatively, maybe I can propose taking 1 potato and 2 papers. That would give me 2 (from potato) + 2 (from 2 papers) = 4 points. Then the other player could take the remaining 1 potato, 2 papers, and 2 clocks. But I don't know their values. If they value clocks more, maybe they want the clocks. But if they take the clocks} \\
            }
        }
    }
    & \\ \\

    \theutterance \stepcounter{utterance}  
    & & & \multicolumn{4}{p{0.6\linewidth}}{
        \cellcolor[rgb]{0.9,0.9,0.9}{
            \makecell[{{p{\linewidth}}}]{
                \texttt{\tiny{[GM$\langle$P2]}}
                \texttt{<think>} \\
\texttt{Okay, let's try to figure out the best way to approach this negotiation. My value function is paper:1, clock:1, potato:2. The other player's values are secret, so I need to make a proposal that's likely to be complementary. The key is to maximize my own score while ensuring that the other player can also get a set of items that doesn't overlap with mine.} \\
\\ 
\texttt{First, let's list all the items: 4 papers, 2 clocks, 2 potatoes. I need to decide how to divide these. Since I value potatoes the most, I should aim to get as many as possible. But there are only two potatoes. If I take both, that's 2*2=4 points. Then the remaining items are 4 papers and 2 clocks. The other player might want some of these. But if they also need potatoes, taking both could cause a conflict. So maybe I should take one potato and leave the other for them. Let's see.} \\
\\ 
\texttt{If I take 1 potato (2 points), then the other player can take the remaining 1 potato. That leaves 4 papers and 2 clocks. If I take some papers or clocks, I can get more points. For example, if I take 1 potato and 2 papers, that's 2 + 2 = 4 points. The other player could take 1 potato, 2 papers, and 2 clocks. But I don't know their values. If they value clocks more, they might prefer the clocks. Alternatively, if they value papers more, they might take the papers.} \\
\\ 
\texttt{Another option: take 1 potato and 1 clock. That's 2 + 1 = 3 points. Then the other player could take 1 potato, 1 clock, and the remaining 4 papers. But maybe that's not optimal for me. I need to maximize my own score while ensuring the other player can get their desired items.} \\
\\ 
\texttt{Wait, the problem says that if the proposals are complementary, meaning there are enough items to fulfill both, then we get our scores. Otherwise, zero. So my proposal must not require more of any item than available, and the other player's proposal must also fit within the remaining items.} \\
\\ 
\texttt{So, if I propose taking 2 potatoes, that's 2 potatoes. The other player's proposal must not require more than 0 potatoes. If they need any potatoes, that would conflict} \\
            }
        }
    }
     \\ \\

    \theutterance \stepcounter{utterance}  
    & & & \multicolumn{2}{p{0.3\linewidth}}{
        \cellcolor[rgb]{0.9,0.9,0.9}{
            \makecell[{{p{\linewidth}}}]{
                \texttt{\tiny{[GM$|$GM]}}
                \texttt{continue} \\
            }
        }
    }
    & & \\ \\

    \theutterance \stepcounter{utterance}  
    & & & \multicolumn{2}{p{0.3\linewidth}}{
        \cellcolor[rgb]{0.9,0.9,0.9}{
            \makecell[{{p{\linewidth}}}]{
                \texttt{\tiny{[GM$|$GM]}}
                \texttt{<think>} \\
\texttt{Okay, let's try to figure out the best way to approach this negotiation. My value function is paper:1, clock:1, potato:2. The other player's values are secret, so I need to make a proposal that's likely to be complementary. The key is to maximize my own score while ensuring that the other player can also get a set of items that doesn't overlap with mine.} \\
\\ 
\texttt{First, let's list all the items: 4 papers, 2 clocks, 2 potatoes. I need to decide how to divide these. Since I value potatoes the most, I should aim to get as many as possible. But there are only two potatoes. If I take both, that's 2*2=4 points. Then the remaining items are 4 papers and 2 clocks. The other player might want some of these. But if they also need potatoes, taking both could cause a conflict. So maybe I should take one potato and leave the other for them. Let's see.} \\
\\ 
\texttt{If I take 1 potato (2 points), then the other player can take the remaining 1 potato. That leaves 4 papers and 2 clocks. If I take some papers or clocks, I can get more points. For example, if I take 1 potato and 2 papers, that's 2 + 2 = 4 points. The other player could take 1 potato, 2 papers, and 2 clocks. But I don't know their values. If they value clocks more, they might prefer the clocks. Alternatively, if they value papers more, they might take the papers.} \\
\\ 
\texttt{Another option: take 1 potato and 1 clock. That's 2 + 1 = 3 points. Then the other player could take 1 potato, 1 clock, and the remaining 4 papers. But maybe that's not optimal for me. I need to maximize my own score while ensuring the other player can get their desired items.} \\
\\ 
\texttt{Wait, the problem says that if the proposals are complementary, meaning there are enough items to fulfill both, then we get our scores. Otherwise, zero. So my proposal must not require more of any item than available, and the other player's proposal must also fit within the remaining items.} \\
\\ 
\texttt{So, if I propose taking 2 potatoes, that's 2 potatoes. The other player's proposal must not require more than 0 potatoes. If they need any potatoes, that would conflict} \\
            }
        }
    }
    & & \\ \\

    \theutterance \stepcounter{utterance}  
    & & \multicolumn{4}{p{0.6\linewidth}}{
        \cellcolor[rgb]{0.9,0.9,0.9}{
            \makecell[{{p{\linewidth}}}]{
                \texttt{\tiny{[P1$\langle$GM]}}
                \texttt{<think>} \\
\texttt{Okay, let's try to figure out the best way to approach this negotiation. My value function is paper:1, clock:1, potato:2. The other player's values are secret, so I need to make a proposal that's likely to be complementary. The key is to maximize my own score while ensuring that the other player can also get a set of items that doesn't overlap with mine.} \\
\\ 
\texttt{First, let's list all the items: 4 papers, 2 clocks, 2 potatoes. I need to decide how to divide these. Since I value potatoes the most, I should aim to get as many as possible. But there are only two potatoes. If I take both, that's 2*2=4 points. Then the remaining items are 4 papers and 2 clocks. The other player might want some of these. But if they also need potatoes, taking both could cause a conflict. So maybe I should take one potato and leave the other for them. Let's see.} \\
\\ 
\texttt{If I take 1 potato (2 points), then the other player can take the remaining 1 potato. That leaves 4 papers and 2 clocks. If I take some papers or clocks, I can get more points. For example, if I take 1 potato and 2 papers, that's 2 + 2 = 4 points. The other player could take 1 potato, 2 papers, and 2 clocks. But I don't know their values. If they value clocks more, they might prefer the clocks. Alternatively, if they value papers more, they might take the papers.} \\
\\ 
\texttt{Another option: take 1 potato and 1 clock. That's 2 + 1 = 3 points. Then the other player could take 1 potato, 1 clock, and the remaining 4 papers. But maybe that's not optimal for me. I need to maximize my own score while ensuring the other player can get their desired items.} \\
\\ 
\texttt{Wait, the problem says that if the proposals are complementary, meaning there are enough items to fulfill both, then we get our scores. Otherwise, zero. So my proposal must not require more of any item than available, and the other player's proposal must also fit within the remaining items.} \\
\\ 
\texttt{So, if I propose taking 2 potatoes, that's 2 potatoes. The other player's proposal must not require more than 0 potatoes. If they need any potatoes, that would conflict} \\
            }
        }
    }
    & \\ \\

    \theutterance \stepcounter{utterance}  
    & \multicolumn{4}{p{0.6\linewidth}}{
        \cellcolor[rgb]{0.9,0.9,0.9}{
            \makecell[{{p{\linewidth}}}]{
                \texttt{\tiny{[P1$\rangle$GM]}}
                \texttt{<think>} \\
\texttt{Okay, let's try to figure out the best way to approach this negotiation. My value function is paper:1, clock:1, potato:2. The other player's values are secret, so I need to make a proposal that's likely to be complementary. The key is to maximize my own score while ensuring that the other player can also get a set of items that doesn't overlap with mine.} \\
\\ 
\texttt{First, let's list all the items: 4 papers, 2 clocks, 2 potatoes. I need to decide how to divide these. Since I value potatoes the most, I should aim to get as many as possible. But there are only two potatoes. If I take both, that's 2*2=4 points. Then the remaining items are 4 papers and 2 clocks. The other player might want some of these. But if they also need potatoes, taking both could cause a conflict. So maybe I should take one potato and leave the other for them. Let's see.} \\
\\ 
\texttt{If I take 1 potato (2 points), then the other player can take the remaining 1 potato. That leaves 4 papers and 2 clocks. If I take some papers or clocks, I can get more points. For example, if I take 1 potato and 2 papers, that's 2 + 2 = 4 points. The other player could take 1 potato, 2 papers, and 2 clocks. But I don't know their values. If they value clocks more, they might prefer the clocks. Alternatively, if they value papers more, they might take the papers.} \\
\\ 
\texttt{Another option: take 1 potato and 1 clock. That's 2 + 1 = 3 points. Then the other player could take 1 potato, 1 clock, and the remaining 4 papers. But maybe that's not optimal for me. I need to maximize my own score while ensuring the other player can get their desired items.} \\
\\ 
\texttt{Wait, the problem says that if the proposals are complementary, meaning there are enough items to fulfill both, then we get our scores. Otherwise, zero. So my proposal must not require more of any item than available, and the other player's proposal must also fit within the remaining items.} \\
\\ 
\texttt{So, if I propose taking 2 potatoes, that's 2 potatoes. The other player's proposal must not require more than 0 potatoes. If they need any potatoes, that would conflict} \\
            }
        }
    }
    & & \\ \\

    \theutterance \stepcounter{utterance}  
    & & & \multicolumn{2}{p{0.3\linewidth}}{
        \cellcolor[rgb]{0.9,0.9,0.9}{
            \makecell[{{p{\linewidth}}}]{
                \texttt{\tiny{[GM$|$GM]}}
                \texttt{continue} \\
            }
        }
    }
    & & \\ \\

    \theutterance \stepcounter{utterance}  
    & & & \multicolumn{2}{p{0.3\linewidth}}{
        \cellcolor[rgb]{0.9,0.9,0.9}{
            \makecell[{{p{\linewidth}}}]{
                \texttt{\tiny{[GM$|$GM]}}
                \texttt{<think>} \\
\texttt{Okay, let's try to figure out the best way to approach this negotiation. My value function is paper:1, clock:1, potato:2. The other player's values are secret, so I need to make a proposal that's likely to be complementary. The key is to maximize my own score while ensuring that the other player can also get a set of items that doesn't overlap with mine.} \\
\\ 
\texttt{First, let's list all the items: 4 papers, 2 clocks, 2 potatoes. I need to decide how to divide these. Since I value potatoes the most, I should aim to get as many as possible. But there are only two potatoes. If I take both, that's 2*2=4 points. Then the remaining items are 4 papers and 2 clocks. The other player might want some of these. But if they also need potatoes, taking both could cause a conflict. So maybe I should take one potato and leave the other for them. Let's see.} \\
\\ 
\texttt{If I take 1 potato (2 points), then the other player can take the remaining 1 potato. That leaves 4 papers and 2 clocks. If I take some papers or clocks, I can get more points. For example, if I take 1 potato and 2 papers, that's 2 + 2 = 4 points. The other player could take 1 potato, 2 papers, and 2 clocks. But I don't know their values. If they value clocks more, they might prefer the clocks. Alternatively, if they value papers more, they might take the papers.} \\
\\ 
\texttt{Another option: take 1 potato and 1 clock. That's 2 + 1 = 3 points. Then the other player could take 1 potato, 1 clock, and the remaining 4 papers. But maybe that's not optimal for me. I need to maximize my own score while ensuring the other player can get their desired items.} \\
\\ 
\texttt{Wait, the problem says that if the proposals are complementary, meaning there are enough items to fulfill both, then we get our scores. Otherwise, zero. So my proposal must not require more of any item than available, and the other player's proposal must also fit within the remaining items.} \\
\\ 
\texttt{So, if I propose taking 2 potatoes, that's 2 potatoes. The other player's proposal must not require more than 0 potatoes. If they need any potatoes, that would conflict} \\
            }
        }
    }
    & & \\ \\

    \theutterance \stepcounter{utterance}  
    & & \multicolumn{4}{p{0.6\linewidth}}{
        \cellcolor[rgb]{0.9,0.9,0.9}{
            \makecell[{{p{\linewidth}}}]{
                \texttt{\tiny{[GM$\rangle$P2]}}
                \texttt{<think>} \\
\texttt{Okay, let's try to figure out the best way to approach this negotiation. My value function is paper:1, clock:1, potato:2. The other player's values are secret, so I need to make a proposal that's likely to be complementary. The key is to maximize my own score while ensuring that the other player can also get a set of items that doesn't overlap with mine.} \\
\\ 
\texttt{First, let's list all the items: 4 papers, 2 clocks, 2 potatoes. I need to decide how to divide these. Since I value potatoes the most, I should aim to get as many as possible. But there are only two potatoes. If I take both, that's 2*2=4 points. Then the remaining items are 4 papers and 2 clocks. The other player might want some of these. But if they also need potatoes, taking both could cause a conflict. So maybe I should take one potato and leave the other for them. Let's see.} \\
\\ 
\texttt{If I take 1 potato (2 points), then the other player can take the remaining 1 potato. That leaves 4 papers and 2 clocks. If I take some papers or clocks, I can get more points. For example, if I take 1 potato and 2 papers, that's 2 + 2 = 4 points. The other player could take 1 potato, 2 papers, and 2 clocks. But I don't know their values. If they value clocks more, they might prefer the clocks. Alternatively, if they value papers more, they might take the papers.} \\
\\ 
\texttt{Another option: take 1 potato and 1 clock. That's 2 + 1 = 3 points. Then the other player could take 1 potato, 1 clock, and the remaining 4 papers. But maybe that's not optimal for me. I need to maximize my own score while ensuring the other player can get their desired items.} \\
\\ 
\texttt{Wait, the problem says that if the proposals are complementary, meaning there are enough items to fulfill both, then we get our scores. Otherwise, zero. So my proposal must not require more of any item than available, and the other player's proposal must also fit within the remaining items.} \\
\\ 
\texttt{So, if I propose taking 2 potatoes, that's 2 potatoes. The other player's proposal must not require more than 0 potatoes. If they need any potatoes, that would conflict} \\
            }
        }
    }
    & \\ \\

    \theutterance \stepcounter{utterance}  
    & & & \multicolumn{4}{p{0.6\linewidth}}{
        \cellcolor[rgb]{0.9,0.9,0.9}{
            \makecell[{{p{\linewidth}}}]{
                \texttt{\tiny{[GM$\langle$P2]}}
                \texttt{<think>} \\
\texttt{Okay, let's try to figure out the best way to approach this negotiation. My value function is paper:1, clock:1, potato:2. The other player's values are secret, so I need to make a proposal that's likely to be complementary. The key is to maximize my own score while ensuring that the other player can also get a set of items that doesn't overlap with mine.} \\
\\ 
\texttt{First, let's list all the items: 4 papers, 2 clocks, 2 potatoes. I need to decide how to divide these. Since I value potatoes the most, I should aim to get as many as possible. But there are only two potatoes. If I take both, that's 2*2=4 points. Then the remaining items are 4 papers and 2 clocks. The other player might want some of these. But if they also need potatoes, taking both could cause a conflict. So maybe I should take one potato and leave the other for them. Let's see.} \\
\\ 
\texttt{If I take 1 potato (2 points), then the other player can take the remaining 1 potato. That leaves 4 papers and 2 clocks. If I take some papers or clocks, I can get more points. For example, if I take 1 potato and 2 papers, that's 2 + 2 = 4 points. The other player could take 1 potato, 2 papers, and 2 clocks. But I don't know their values. If they value clocks more, they might prefer the clocks. Alternatively, if they value papers more, they might take the papers.} \\
\\ 
\texttt{Another option: take 1 potato and 1 clock. That's 2 + 1 = 3 points. Then the other player could take 1 potato, 1 clock, and the remaining 4 papers. But maybe that's not optimal for me. I need to maximize my own score while ensuring the other player can get their desired items.} \\
\\ 
\texttt{Wait, the problem says that if the proposals are complementary, meaning there are enough items to fulfill both, then we get our scores. Otherwise, zero. So my proposal must not require more of any item than available, and the other player's proposal must also fit within the remaining items.} \\
\\ 
\texttt{So, if I propose taking 2 potatoes, that's 2 potatoes. The other player's proposal must not require more than 0 potatoes. If they need any potatoes, that would conflict} \\
            }
        }
    }
     \\ \\

    \theutterance \stepcounter{utterance}  
    & & & \multicolumn{2}{p{0.3\linewidth}}{
        \cellcolor[rgb]{0.9,0.9,0.9}{
            \makecell[{{p{\linewidth}}}]{
                \texttt{\tiny{[GM$|$GM]}}
                \texttt{continue} \\
            }
        }
    }
    & & \\ \\

    \theutterance \stepcounter{utterance}  
    & & & \multicolumn{2}{p{0.3\linewidth}}{
        \cellcolor[rgb]{0.9,0.9,0.9}{
            \makecell[{{p{\linewidth}}}]{
                \texttt{\tiny{[GM$|$GM]}}
                \texttt{<think>} \\
\texttt{Okay, let's try to figure out the best way to approach this negotiation. My value function is paper:1, clock:1, potato:2. The other player's values are secret, so I need to make a proposal that's likely to be complementary. The key is to maximize my own score while ensuring that the other player can also get a set of items that doesn't overlap with mine.} \\
\\ 
\texttt{First, let's list all the items: 4 papers, 2 clocks, 2 potatoes. I need to decide how to divide these. Since I value potatoes the most, I should aim to get as many as possible. But there are only two potatoes. If I take both, that's 2*2=4 points. Then the remaining items are 4 papers and 2 clocks. The other player might want some of these. But if they also need potatoes, taking both could cause a conflict. So maybe I should take one potato and leave the other for them. Let's see.} \\
\\ 
\texttt{If I take 1 potato (2 points), then the other player can take the remaining 1 potato. That leaves 4 papers and 2 clocks. If I take some papers or clocks, I can get more points. For example, if I take 1 potato and 2 papers, that's 2 + 2 = 4 points. The other player could take 1 potato, 2 papers, and 2 clocks. But I don't know their values. If they value clocks more, they might prefer the clocks. Alternatively, if they value papers more, they might take the papers.} \\
\\ 
\texttt{Another option: take 1 potato and 1 clock. That's 2 + 1 = 3 points. Then the other player could take 1 potato, 1 clock, and the remaining 4 papers. But maybe that's not optimal for me. I need to maximize my own score while ensuring the other player can get their desired items.} \\
\\ 
\texttt{Wait, the problem says that if the proposals are complementary, meaning there are enough items to fulfill both, then we get our scores. Otherwise, zero. So my proposal must not require more of any item than available, and the other player's proposal must also fit within the remaining items.} \\
\\ 
\texttt{So, if I propose taking 2 potatoes, that's 2 potatoes. The other player's proposal must not require more than 0 potatoes. If they need any potatoes, that would conflict} \\
            }
        }
    }
    & & \\ \\

    \theutterance \stepcounter{utterance}  
    & & \multicolumn{4}{p{0.6\linewidth}}{
        \cellcolor[rgb]{0.9,0.9,0.9}{
            \makecell[{{p{\linewidth}}}]{
                \texttt{\tiny{[P1$\langle$GM]}}
                \texttt{<think>} \\
\texttt{Okay, let's try to figure out the best way to approach this negotiation. My value function is paper:1, clock:1, potato:2. The other player's values are secret, so I need to make a proposal that's likely to be complementary. The key is to maximize my own score while ensuring that the other player can also get a set of items that doesn't overlap with mine.} \\
\\ 
\texttt{First, let's list all the items: 4 papers, 2 clocks, 2 potatoes. I need to decide how to divide these. Since I value potatoes the most, I should aim to get as many as possible. But there are only two potatoes. If I take both, that's 2*2=4 points. Then the remaining items are 4 papers and 2 clocks. The other player might want some of these. But if they also need potatoes, taking both could cause a conflict. So maybe I should take one potato and leave the other for them. Let's see.} \\
\\ 
\texttt{If I take 1 potato (2 points), then the other player can take the remaining 1 potato. That leaves 4 papers and 2 clocks. If I take some papers or clocks, I can get more points. For example, if I take 1 potato and 2 papers, that's 2 + 2 = 4 points. The other player could take 1 potato, 2 papers, and 2 clocks. But I don't know their values. If they value clocks more, they might prefer the clocks. Alternatively, if they value papers more, they might take the papers.} \\
\\ 
\texttt{Another option: take 1 potato and 1 clock. That's 2 + 1 = 3 points. Then the other player could take 1 potato, 1 clock, and the remaining 4 papers. But maybe that's not optimal for me. I need to maximize my own score while ensuring the other player can get their desired items.} \\
\\ 
\texttt{Wait, the problem says that if the proposals are complementary, meaning there are enough items to fulfill both, then we get our scores. Otherwise, zero. So my proposal must not require more of any item than available, and the other player's proposal must also fit within the remaining items.} \\
\\ 
\texttt{So, if I propose taking 2 potatoes, that's 2 potatoes. The other player's proposal must not require more than 0 potatoes. If they need any potatoes, that would conflict} \\
            }
        }
    }
    & \\ \\

    \theutterance \stepcounter{utterance}  
    & \multicolumn{4}{p{0.6\linewidth}}{
        \cellcolor[rgb]{0.9,0.9,0.9}{
            \makecell[{{p{\linewidth}}}]{
                \texttt{\tiny{[P1$\rangle$GM]}}
                \texttt{<think>} \\
\texttt{Okay, let's try to figure out the best way to approach this negotiation. My value function is paper:1, clock:1, potato:2. The other player's values are secret, so I need to make a proposal that's likely to be complementary. The key is to maximize my own score while ensuring that the other player can also get a set of items that doesn't overlap with mine.} \\
\\ 
\texttt{First, let's list all the items: 4 papers, 2 clocks, 2 potatoes. I need to decide how to divide these. Since I value potatoes the most, I should aim to get as many as possible. But there are only two potatoes. If I take both, that's 2*2=4 points. Then the remaining items are 4 papers and 2 clocks. The other player might want some of these. But if they also need potatoes, taking both could cause a conflict. So maybe I should take one potato and leave the other for them. Let's see.} \\
\\ 
\texttt{If I take 1 potato (2 points), then the other player can take the remaining 1 potato. That leaves 4 papers and 2 clocks. If I take some papers or clocks, I can get more points. For example, if I take 1 potato and 2 papers, that's 2 + 2 = 4 points. The other player could take 1 potato, 2 papers, and 2 clocks. But I don't know their values. If they value clocks more, they might prefer the clocks. Alternatively, if they value papers more, they might take the papers.} \\
\\ 
\texttt{Another option: take 1 potato and 1 clock. That's 2 + 1 = 3 points. Then the other player could take 1 potato, 1 clock, and the remaining 4 papers. But maybe that's not optimal for me. I need to maximize my own score while ensuring the other player can get their desired items.} \\
\\ 
\texttt{Wait, the problem says that if the proposals are complementary, meaning there are enough items to fulfill both, then we get our scores. Otherwise, zero. So my proposal must not require more of any item than available, and the other player's proposal must also fit within the remaining items.} \\
\\ 
\texttt{So, if I propose taking 2 potatoes, that's 2 potatoes. The other player's proposal must not require more than 0 potatoes. If they need any potatoes, that would conflict} \\
            }
        }
    }
    & & \\ \\

    \theutterance \stepcounter{utterance}  
    & & & \multicolumn{2}{p{0.3\linewidth}}{
        \cellcolor[rgb]{0.9,0.9,0.9}{
            \makecell[{{p{\linewidth}}}]{
                \texttt{\tiny{[GM$|$GM]}}
                \texttt{continue} \\
            }
        }
    }
    & & \\ \\

    \theutterance \stepcounter{utterance}  
    & & & \multicolumn{2}{p{0.3\linewidth}}{
        \cellcolor[rgb]{0.9,0.9,0.9}{
            \makecell[{{p{\linewidth}}}]{
                \texttt{\tiny{[GM$|$GM]}}
                \texttt{<think>} \\
\texttt{Okay, let's try to figure out the best way to approach this negotiation. My value function is paper:1, clock:1, potato:2. The other player's values are secret, so I need to make a proposal that's likely to be complementary. The key is to maximize my own score while ensuring that the other player can also get a set of items that doesn't overlap with mine.} \\
\\ 
\texttt{First, let's list all the items: 4 papers, 2 clocks, 2 potatoes. I need to decide how to divide these. Since I value potatoes the most, I should aim to get as many as possible. But there are only two potatoes. If I take both, that's 2*2=4 points. Then the remaining items are 4 papers and 2 clocks. The other player might want some of these. But if they also need potatoes, taking both could cause a conflict. So maybe I should take one potato and leave the other for them. Let's see.} \\
\\ 
\texttt{If I take 1 potato (2 points), then the other player can take the remaining 1 potato. That leaves 4 papers and 2 clocks. If I take some papers or clocks, I can get more points. For example, if I take 1 potato and 2 papers, that's 2 + 2 = 4 points. The other player could take 1 potato, 2 papers, and 2 clocks. But I don't know their values. If they value clocks more, they might prefer the clocks. Alternatively, if they value papers more, they might take the papers.} \\
\\ 
\texttt{Another option: take 1 potato and 1 clock. That's 2 + 1 = 3 points. Then the other player could take 1 potato, 1 clock, and the remaining 4 papers. But maybe that's not optimal for me. I need to maximize my own score while ensuring the other player can get their desired items.} \\
\\ 
\texttt{Wait, the problem says that if the proposals are complementary, meaning there are enough items to fulfill both, then we get our scores. Otherwise, zero. So my proposal must not require more of any item than available, and the other player's proposal must also fit within the remaining items.} \\
\\ 
\texttt{So, if I propose taking 2 potatoes, that's 2 potatoes. The other player's proposal must not require more than 0 potatoes. If they need any potatoes, that would conflict} \\
            }
        }
    }
    & & \\ \\

    \theutterance \stepcounter{utterance}  
    & & \multicolumn{4}{p{0.6\linewidth}}{
        \cellcolor[rgb]{0.9,0.9,0.9}{
            \makecell[{{p{\linewidth}}}]{
                \texttt{\tiny{[GM$\rangle$P2]}}
                \texttt{<think>} \\
\texttt{Okay, let's try to figure out the best way to approach this negotiation. My value function is paper:1, clock:1, potato:2. The other player's values are secret, so I need to make a proposal that's likely to be complementary. The key is to maximize my own score while ensuring that the other player can also get a set of items that doesn't overlap with mine.} \\
\\ 
\texttt{First, let's list all the items: 4 papers, 2 clocks, 2 potatoes. I need to decide how to divide these. Since I value potatoes the most, I should aim to get as many as possible. But there are only two potatoes. If I take both, that's 2*2=4 points. Then the remaining items are 4 papers and 2 clocks. The other player might want some of these. But if they also need potatoes, taking both could cause a conflict. So maybe I should take one potato and leave the other for them. Let's see.} \\
\\ 
\texttt{If I take 1 potato (2 points), then the other player can take the remaining 1 potato. That leaves 4 papers and 2 clocks. If I take some papers or clocks, I can get more points. For example, if I take 1 potato and 2 papers, that's 2 + 2 = 4 points. The other player could take 1 potato, 2 papers, and 2 clocks. But I don't know their values. If they value clocks more, they might prefer the clocks. Alternatively, if they value papers more, they might take the papers.} \\
\\ 
\texttt{Another option: take 1 potato and 1 clock. That's 2 + 1 = 3 points. Then the other player could take 1 potato, 1 clock, and the remaining 4 papers. But maybe that's not optimal for me. I need to maximize my own score while ensuring the other player can get their desired items.} \\
\\ 
\texttt{Wait, the problem says that if the proposals are complementary, meaning there are enough items to fulfill both, then we get our scores. Otherwise, zero. So my proposal must not require more of any item than available, and the other player's proposal must also fit within the remaining items.} \\
\\ 
\texttt{So, if I propose taking 2 potatoes, that's 2 potatoes. The other player's proposal must not require more than 0 potatoes. If they need any potatoes, that would conflict} \\
            }
        }
    }
    & \\ \\

    \theutterance \stepcounter{utterance}  
    & & & \multicolumn{4}{p{0.6\linewidth}}{
        \cellcolor[rgb]{0.9,0.9,0.9}{
            \makecell[{{p{\linewidth}}}]{
                \texttt{\tiny{[GM$\langle$P2]}}
                \texttt{<think>} \\
\texttt{Okay, let's try to figure out the best way to approach this negotiation. My value function is paper:1, clock:1, potato:2. The other player's values are secret, so I need to make a proposal that's likely to be complementary. The key is to maximize my own score while ensuring that the other player can also get a set of items that doesn't overlap with mine.} \\
\\ 
\texttt{First, let's list all the items: 4 papers, 2 clocks, 2 potatoes. I need to decide how to divide these. Since I value potatoes the most, I should aim to get as many as possible. But there are only two potatoes. If I take both, that's 2*2=4 points. Then the remaining items are 4 papers and 2 clocks. The other player might want some of these. But if they also need potatoes, taking both could cause a conflict. So maybe I should take one potato and leave the other for them. Let's see.} \\
\\ 
\texttt{If I take 1 potato (2 points), then the other player can take the remaining 1 potato. That leaves 4 papers and 2 clocks. If I take some papers or clocks, I can get more points. For example, if I take 1 potato and 2 papers, that's 2 + 2 = 4 points. The other player could take 1 potato, 2 papers, and 2 clocks. But I don't know their values. If they value clocks more, they might prefer the clocks. Alternatively, if they value papers more, they might take the papers.} \\
\\ 
\texttt{Another option: take 1 potato and 1 clock. That's 2 + 1 = 3 points. Then the other player could take 1 potato, 1 clock, and the remaining 4 papers. But maybe that's not optimal for me. I need to maximize my own score while ensuring the other player can get their desired items.} \\
\\ 
\texttt{Wait, the problem says that if the proposals are complementary, meaning there are enough items to fulfill both, then we get our scores. Otherwise, zero. So my proposal must not require more of any item than available, and the other player's proposal must also fit within the remaining items.} \\
\\ 
\texttt{So, if I propose taking 2 potatoes, that's 2 potatoes. The other player's proposal must not require more than 0 potatoes. If they need any potatoes, that would conflict} \\
            }
        }
    }
     \\ \\

    \theutterance \stepcounter{utterance}  
    & & & \multicolumn{2}{p{0.3\linewidth}}{
        \cellcolor[rgb]{0.9,0.9,0.9}{
            \makecell[{{p{\linewidth}}}]{
                \texttt{\tiny{[GM$|$GM]}}
                \texttt{continue} \\
            }
        }
    }
    & & \\ \\

    \theutterance \stepcounter{utterance}  
    & & & \multicolumn{2}{p{0.3\linewidth}}{
        \cellcolor[rgb]{0.9,0.9,0.9}{
            \makecell[{{p{\linewidth}}}]{
                \texttt{\tiny{[GM$|$GM]}}
                \texttt{<think>} \\
\texttt{Okay, let's try to figure out the best way to approach this negotiation. My value function is paper:1, clock:1, potato:2. The other player's values are secret, so I need to make a proposal that's likely to be complementary. The key is to maximize my own score while ensuring that the other player can also get a set of items that doesn't overlap with mine.} \\
\\ 
\texttt{First, let's list all the items: 4 papers, 2 clocks, 2 potatoes. I need to decide how to divide these. Since I value potatoes the most, I should aim to get as many as possible. But there are only two potatoes. If I take both, that's 2*2=4 points. Then the remaining items are 4 papers and 2 clocks. The other player might want some of these. But if they also need potatoes, taking both could cause a conflict. So maybe I should take one potato and leave the other for them. Let's see.} \\
\\ 
\texttt{If I take 1 potato (2 points), then the other player can take the remaining 1 potato. That leaves 4 papers and 2 clocks. If I take some papers or clocks, I can get more points. For example, if I take 1 potato and 2 papers, that's 2 + 2 = 4 points. The other player could take 1 potato, 2 papers, and 2 clocks. But I don't know their values. If they value clocks more, they might prefer the clocks. Alternatively, if they value papers more, they might take the papers.} \\
\\ 
\texttt{Another option: take 1 potato and 1 clock. That's 2 + 1 = 3 points. Then the other player could take 1 potato, 1 clock, and the remaining 4 papers. But maybe that's not optimal for me. I need to maximize my own score while ensuring the other player can get their desired items.} \\
\\ 
\texttt{Wait, the problem says that if the proposals are complementary, meaning there are enough items to fulfill both, then we get our scores. Otherwise, zero. So my proposal must not require more of any item than available, and the other player's proposal must also fit within the remaining items.} \\
\\ 
\texttt{So, if I propose taking 2 potatoes, that's 2 potatoes. The other player's proposal must not require more than 0 potatoes. If they need any potatoes, that would conflict} \\
            }
        }
    }
    & & \\ \\

    \theutterance \stepcounter{utterance}  
    & & \multicolumn{4}{p{0.6\linewidth}}{
        \cellcolor[rgb]{0.9,0.9,0.9}{
            \makecell[{{p{\linewidth}}}]{
                \texttt{\tiny{[P1$\langle$GM]}}
                \texttt{<think>} \\
\texttt{Okay, let's try to figure out the best way to approach this negotiation. My value function is paper:1, clock:1, potato:2. The other player's values are secret, so I need to make a proposal that's likely to be complementary. The key is to maximize my own score while ensuring that the other player can also get a set of items that doesn't overlap with mine.} \\
\\ 
\texttt{First, let's list all the items: 4 papers, 2 clocks, 2 potatoes. I need to decide how to divide these. Since I value potatoes the most, I should aim to get as many as possible. But there are only two potatoes. If I take both, that's 2*2=4 points. Then the remaining items are 4 papers and 2 clocks. The other player might want some of these. But if they also need potatoes, taking both could cause a conflict. So maybe I should take one potato and leave the other for them. Let's see.} \\
\\ 
\texttt{If I take 1 potato (2 points), then the other player can take the remaining 1 potato. That leaves 4 papers and 2 clocks. If I take some papers or clocks, I can get more points. For example, if I take 1 potato and 2 papers, that's 2 + 2 = 4 points. The other player could take 1 potato, 2 papers, and 2 clocks. But I don't know their values. If they value clocks more, they might prefer the clocks. Alternatively, if they value papers more, they might take the papers.} \\
\\ 
\texttt{Another option: take 1 potato and 1 clock. That's 2 + 1 = 3 points. Then the other player could take 1 potato, 1 clock, and the remaining 4 papers. But maybe that's not optimal for me. I need to maximize my own score while ensuring the other player can get their desired items.} \\
\\ 
\texttt{Wait, the problem says that if the proposals are complementary, meaning there are enough items to fulfill both, then we get our scores. Otherwise, zero. So my proposal must not require more of any item than available, and the other player's proposal must also fit within the remaining items.} \\
\\ 
\texttt{So, if I propose taking 2 potatoes, that's 2 potatoes. The other player's proposal must not require more than 0 potatoes. If they need any potatoes, that would conflict} \\
            }
        }
    }
    & \\ \\

    \theutterance \stepcounter{utterance}  
    & \multicolumn{4}{p{0.6\linewidth}}{
        \cellcolor[rgb]{0.9,0.9,0.9}{
            \makecell[{{p{\linewidth}}}]{
                \texttt{\tiny{[P1$\rangle$GM]}}
                \texttt{<think>} \\
\texttt{Okay, let's try to figure out the best way to approach this negotiation. My value function is paper:1, clock:1, potato:2. The other player's values are secret, so I need to make a proposal that's likely to be complementary. The key is to maximize my own score while ensuring that the other player can also get a set of items that doesn't overlap with mine.} \\
\\ 
\texttt{First, let's list all the items: 4 papers, 2 clocks, 2 potatoes. I need to decide how to divide these. Since I value potatoes the most, I should aim to get as many as possible. But there are only two potatoes. If I take both, that's 2*2=4 points. Then the remaining items are 4 papers and 2 clocks. The other player might want some of these. But if they also need potatoes, taking both could cause a conflict. So maybe I should take one potato and leave the other for them. Let's see.} \\
\\ 
\texttt{If I take 1 potato (2 points), then the other player can take the remaining 1 potato. That leaves 4 papers and 2 clocks. If I take some papers or clocks, I can get more points. For example, if I take 1 potato and 2 papers, that's 2 + 2 = 4 points. The other player could take 1 potato, 2 papers, and 2 clocks. But I don't know their values. If they value clocks more, they might prefer the clocks. Alternatively, if they value papers more, they might take the papers.} \\
\\ 
\texttt{Another option: take 1 potato and 1 clock. That's 2 + 1 = 3 points. Then the other player could take 1 potato, 1 clock, and the remaining 4 papers. But maybe that's not optimal for me. I need to maximize my own score while ensuring the other player can get their desired items.} \\
\\ 
\texttt{Wait, the problem says that if the proposals are complementary, meaning there are enough items to fulfill both, then we get our scores. Otherwise, zero. So my proposal must not require more of any item than available, and the other player's proposal must also fit within the remaining items.} \\
\\ 
\texttt{So, if I propose taking 2 potatoes, that's 2 potatoes. The other player's proposal must not require more than 0 potatoes. If they need any potatoes, that would conflict} \\
            }
        }
    }
    & & \\ \\

    \theutterance \stepcounter{utterance}  
    & & & \multicolumn{2}{p{0.3\linewidth}}{
        \cellcolor[rgb]{0.9,0.9,0.9}{
            \makecell[{{p{\linewidth}}}]{
                \texttt{\tiny{[GM$|$GM]}}
                \texttt{continue} \\
            }
        }
    }
    & & \\ \\

    \theutterance \stepcounter{utterance}  
    & & & \multicolumn{2}{p{0.3\linewidth}}{
        \cellcolor[rgb]{0.9,0.9,0.9}{
            \makecell[{{p{\linewidth}}}]{
                \texttt{\tiny{[GM$|$GM]}}
                \texttt{<think>} \\
\texttt{Okay, let's try to figure out the best way to approach this negotiation. My value function is paper:1, clock:1, potato:2. The other player's values are secret, so I need to make a proposal that's likely to be complementary. The key is to maximize my own score while ensuring that the other player can also get a set of items that doesn't overlap with mine.} \\
\\ 
\texttt{First, let's list all the items: 4 papers, 2 clocks, 2 potatoes. I need to decide how to divide these. Since I value potatoes the most, I should aim to get as many as possible. But there are only two potatoes. If I take both, that's 2*2=4 points. Then the remaining items are 4 papers and 2 clocks. The other player might want some of these. But if they also need potatoes, taking both could cause a conflict. So maybe I should take one potato and leave the other for them. Let's see.} \\
\\ 
\texttt{If I take 1 potato (2 points), then the other player can take the remaining 1 potato. That leaves 4 papers and 2 clocks. If I take some papers or clocks, I can get more points. For example, if I take 1 potato and 2 papers, that's 2 + 2 = 4 points. The other player could take 1 potato, 2 papers, and 2 clocks. But I don't know their values. If they value clocks more, they might prefer the clocks. Alternatively, if they value papers more, they might take the papers.} \\
\\ 
\texttt{Another option: take 1 potato and 1 clock. That's 2 + 1 = 3 points. Then the other player could take 1 potato, 1 clock, and the remaining 4 papers. But maybe that's not optimal for me. I need to maximize my own score while ensuring the other player can get their desired items.} \\
\\ 
\texttt{Wait, the problem says that if the proposals are complementary, meaning there are enough items to fulfill both, then we get our scores. Otherwise, zero. So my proposal must not require more of any item than available, and the other player's proposal must also fit within the remaining items.} \\
\\ 
\texttt{So, if I propose taking 2 potatoes, that's 2 potatoes. The other player's proposal must not require more than 0 potatoes. If they need any potatoes, that would conflict} \\
            }
        }
    }
    & & \\ \\

    \theutterance \stepcounter{utterance}  
    & & \multicolumn{4}{p{0.6\linewidth}}{
        \cellcolor[rgb]{0.9,0.9,0.9}{
            \makecell[{{p{\linewidth}}}]{
                \texttt{\tiny{[GM$\rangle$P2]}}
                \texttt{<think>} \\
\texttt{Okay, let's try to figure out the best way to approach this negotiation. My value function is paper:1, clock:1, potato:2. The other player's values are secret, so I need to make a proposal that's likely to be complementary. The key is to maximize my own score while ensuring that the other player can also get a set of items that doesn't overlap with mine.} \\
\\ 
\texttt{First, let's list all the items: 4 papers, 2 clocks, 2 potatoes. I need to decide how to divide these. Since I value potatoes the most, I should aim to get as many as possible. But there are only two potatoes. If I take both, that's 2*2=4 points. Then the remaining items are 4 papers and 2 clocks. The other player might want some of these. But if they also need potatoes, taking both could cause a conflict. So maybe I should take one potato and leave the other for them. Let's see.} \\
\\ 
\texttt{If I take 1 potato (2 points), then the other player can take the remaining 1 potato. That leaves 4 papers and 2 clocks. If I take some papers or clocks, I can get more points. For example, if I take 1 potato and 2 papers, that's 2 + 2 = 4 points. The other player could take 1 potato, 2 papers, and 2 clocks. But I don't know their values. If they value clocks more, they might prefer the clocks. Alternatively, if they value papers more, they might take the papers.} \\
\\ 
\texttt{Another option: take 1 potato and 1 clock. That's 2 + 1 = 3 points. Then the other player could take 1 potato, 1 clock, and the remaining 4 papers. But maybe that's not optimal for me. I need to maximize my own score while ensuring the other player can get their desired items.} \\
\\ 
\texttt{Wait, the problem says that if the proposals are complementary, meaning there are enough items to fulfill both, then we get our scores. Otherwise, zero. So my proposal must not require more of any item than available, and the other player's proposal must also fit within the remaining items.} \\
\\ 
\texttt{So, if I propose taking 2 potatoes, that's 2 potatoes. The other player's proposal must not require more than 0 potatoes. If they need any potatoes, that would conflict} \\
            }
        }
    }
    & \\ \\

    \theutterance \stepcounter{utterance}  
    & & & \multicolumn{4}{p{0.6\linewidth}}{
        \cellcolor[rgb]{0.9,0.9,0.9}{
            \makecell[{{p{\linewidth}}}]{
                \texttt{\tiny{[GM$\langle$P2]}}
                \texttt{<think>} \\
\texttt{Okay, let's try to figure out the best way to approach this negotiation. My value function is paper:1, clock:1, potato:2. The other player's values are secret, so I need to make a proposal that's likely to be complementary. The key is to maximize my own score while ensuring that the other player can also get a set of items that doesn't overlap with mine.} \\
\\ 
\texttt{First, let's list all the items: 4 papers, 2 clocks, 2 potatoes. I need to decide how to divide these. Since I value potatoes the most, I should aim to get as many as possible. But there are only two potatoes. If I take both, that's 2*2=4 points. Then the remaining items are 4 papers and 2 clocks. The other player might want some of these. But if they also need potatoes, taking both could cause a conflict. So maybe I should take one potato and leave the other for them. Let's see.} \\
\\ 
\texttt{If I take 1 potato (2 points), then the other player can take the remaining 1 potato. That leaves 4 papers and 2 clocks. If I take some papers or clocks, I can get more points. For example, if I take 1 potato and 2 papers, that's 2 + 2 = 4 points. The other player could take 1 potato, 2 papers, and 2 clocks. But I don't know their values. If they value clocks more, they might prefer the clocks. Alternatively, if they value papers more, they might take the papers.} \\
\\ 
\texttt{Another option: take 1 potato and 1 clock. That's 2 + 1 = 3 points. Then the other player could take 1 potato, 1 clock, and the remaining 4 papers. But maybe that's not optimal for me. I need to maximize my own score while ensuring the other player can get their desired items.} \\
\\ 
\texttt{Wait, the problem says that if the proposals are complementary, meaning there are enough items to fulfill both, then we get our scores. Otherwise, zero. So my proposal must not require more of any item than available, and the other player's proposal must also fit within the remaining items.} \\
\\ 
\texttt{So, if I propose taking 2 potatoes, that's 2 potatoes. The other player's proposal must not require more than 0 potatoes. If they need any potatoes, that would conflict} \\
            }
        }
    }
     \\ \\

    \theutterance \stepcounter{utterance}  
    & & & \multicolumn{2}{p{0.3\linewidth}}{
        \cellcolor[rgb]{0.9,0.9,0.9}{
            \makecell[{{p{\linewidth}}}]{
                \texttt{\tiny{[GM$|$GM]}}
                \texttt{continue} \\
            }
        }
    }
    & & \\ \\

    \theutterance \stepcounter{utterance}  
    & & & \multicolumn{2}{p{0.3\linewidth}}{
        \cellcolor[rgb]{0.9,0.9,0.9}{
            \makecell[{{p{\linewidth}}}]{
                \texttt{\tiny{[GM$|$GM]}}
                \texttt{<think>} \\
\texttt{Okay, let's try to figure out the best way to approach this negotiation. My value function is paper:1, clock:1, potato:2. The other player's values are secret, so I need to make a proposal that's likely to be complementary. The key is to maximize my own score while ensuring that the other player can also get a set of items that doesn't overlap with mine.} \\
\\ 
\texttt{First, let's list all the items: 4 papers, 2 clocks, 2 potatoes. I need to decide how to divide these. Since I value potatoes the most, I should aim to get as many as possible. But there are only two potatoes. If I take both, that's 2*2=4 points. Then the remaining items are 4 papers and 2 clocks. The other player might want some of these. But if they also need potatoes, taking both could cause a conflict. So maybe I should take one potato and leave the other for them. Let's see.} \\
\\ 
\texttt{If I take 1 potato (2 points), then the other player can take the remaining 1 potato. That leaves 4 papers and 2 clocks. If I take some papers or clocks, I can get more points. For example, if I take 1 potato and 2 papers, that's 2 + 2 = 4 points. The other player could take 1 potato, 2 papers, and 2 clocks. But I don't know their values. If they value clocks more, they might prefer the clocks. Alternatively, if they value papers more, they might take the papers.} \\
\\ 
\texttt{Another option: take 1 potato and 1 clock. That's 2 + 1 = 3 points. Then the other player could take 1 potato, 1 clock, and the remaining 4 papers. But maybe that's not optimal for me. I need to maximize my own score while ensuring the other player can get their desired items.} \\
\\ 
\texttt{Wait, the problem says that if the proposals are complementary, meaning there are enough items to fulfill both, then we get our scores. Otherwise, zero. So my proposal must not require more of any item than available, and the other player's proposal must also fit within the remaining items.} \\
\\ 
\texttt{So, if I propose taking 2 potatoes, that's 2 potatoes. The other player's proposal must not require more than 0 potatoes. If they need any potatoes, that would conflict} \\
            }
        }
    }
    & & \\ \\

    \theutterance \stepcounter{utterance}  
    & & \multicolumn{4}{p{0.6\linewidth}}{
        \cellcolor[rgb]{0.9,0.9,0.9}{
            \makecell[{{p{\linewidth}}}]{
                \texttt{\tiny{[P1$\langle$GM]}}
                \texttt{<think>} \\
\texttt{Okay, let's try to figure out the best way to approach this negotiation. My value function is paper:1, clock:1, potato:2. The other player's values are secret, so I need to make a proposal that's likely to be complementary. The key is to maximize my own score while ensuring that the other player can also get a set of items that doesn't overlap with mine.} \\
\\ 
\texttt{First, let's list all the items: 4 papers, 2 clocks, 2 potatoes. I need to decide how to divide these. Since I value potatoes the most, I should aim to get as many as possible. But there are only two potatoes. If I take both, that's 2*2=4 points. Then the remaining items are 4 papers and 2 clocks. The other player might want some of these. But if they also need potatoes, taking both could cause a conflict. So maybe I should take one potato and leave the other for them. Let's see.} \\
\\ 
\texttt{If I take 1 potato (2 points), then the other player can take the remaining 1 potato. That leaves 4 papers and 2 clocks. If I take some papers or clocks, I can get more points. For example, if I take 1 potato and 2 papers, that's 2 + 2 = 4 points. The other player could take 1 potato, 2 papers, and 2 clocks. But I don't know their values. If they value clocks more, they might prefer the clocks. Alternatively, if they value papers more, they might take the papers.} \\
\\ 
\texttt{Another option: take 1 potato and 1 clock. That's 2 + 1 = 3 points. Then the other player could take 1 potato, 1 clock, and the remaining 4 papers. But maybe that's not optimal for me. I need to maximize my own score while ensuring the other player can get their desired items.} \\
\\ 
\texttt{Wait, the problem says that if the proposals are complementary, meaning there are enough items to fulfill both, then we get our scores. Otherwise, zero. So my proposal must not require more of any item than available, and the other player's proposal must also fit within the remaining items.} \\
\\ 
\texttt{So, if I propose taking 2 potatoes, that's 2 potatoes. The other player's proposal must not require more than 0 potatoes. If they need any potatoes, that would conflict} \\
            }
        }
    }
    & \\ \\

    \theutterance \stepcounter{utterance}  
    & \multicolumn{4}{p{0.6\linewidth}}{
        \cellcolor[rgb]{0.9,0.9,0.9}{
            \makecell[{{p{\linewidth}}}]{
                \texttt{\tiny{[P1$\rangle$GM]}}
                \texttt{<think>} \\
\texttt{Okay, let's try to figure out the best way to approach this negotiation. My value function is paper:1, clock:1, potato:2. The other player's values are secret, so I need to make a proposal that's likely to be complementary. The key is to maximize my own score while ensuring that the other player can also get a set of items that doesn't overlap with mine.} \\
\\ 
\texttt{First, let's list all the items: 4 papers, 2 clocks, 2 potatoes. I need to decide how to divide these. Since I value potatoes the most, I should aim to get as many as possible. But there are only two potatoes. If I take both, that's 2*2=4 points. Then the remaining items are 4 papers and 2 clocks. The other player might want some of these. But if they also need potatoes, taking both could cause a conflict. So maybe I should take one potato and leave the other for them. Let's see.} \\
\\ 
\texttt{If I take 1 potato (2 points), then the other player can take the remaining 1 potato. That leaves 4 papers and 2 clocks. If I take some papers or clocks, I can get more points. For example, if I take 1 potato and 2 papers, that's 2 + 2 = 4 points. The other player could take 1 potato, 2 papers, and 2 clocks. But I don't know their values. If they value clocks more, they might prefer the clocks. Alternatively, if they value papers more, they might take the papers.} \\
\\ 
\texttt{Another option: take 1 potato and 1 clock. That's 2 + 1 = 3 points. Then the other player could take 1 potato, 1 clock, and the remaining 4 papers. But maybe that's not optimal for me. I need to maximize my own score while ensuring the other player can get their desired items.} \\
\\ 
\texttt{Wait, the problem says that if the proposals are complementary, meaning there are enough items to fulfill both, then we get our scores. Otherwise, zero. So my proposal must not require more of any item than available, and the other player's proposal must also fit within the remaining items.} \\
\\ 
\texttt{So, if I propose taking 2 potatoes, that's 2 potatoes. The other player's proposal must not require more than 0 potatoes. If they need any potatoes, that would conflict} \\
            }
        }
    }
    & & \\ \\

    \theutterance \stepcounter{utterance}  
    & & & \multicolumn{2}{p{0.3\linewidth}}{
        \cellcolor[rgb]{0.9,0.9,0.9}{
            \makecell[{{p{\linewidth}}}]{
                \texttt{\tiny{[GM$|$GM]}}
                \texttt{continue} \\
            }
        }
    }
    & & \\ \\

    \theutterance \stepcounter{utterance}  
    & & & \multicolumn{2}{p{0.3\linewidth}}{
        \cellcolor[rgb]{0.9,0.9,0.9}{
            \makecell[{{p{\linewidth}}}]{
                \texttt{\tiny{[GM$|$GM]}}
                \texttt{<think>} \\
\texttt{Okay, let's try to figure out the best way to approach this negotiation. My value function is paper:1, clock:1, potato:2. The other player's values are secret, so I need to make a proposal that's likely to be complementary. The key is to maximize my own score while ensuring that the other player can also get a set of items that doesn't overlap with mine.} \\
\\ 
\texttt{First, let's list all the items: 4 papers, 2 clocks, 2 potatoes. I need to decide how to divide these. Since I value potatoes the most, I should aim to get as many as possible. But there are only two potatoes. If I take both, that's 2*2=4 points. Then the remaining items are 4 papers and 2 clocks. The other player might want some of these. But if they also need potatoes, taking both could cause a conflict. So maybe I should take one potato and leave the other for them. Let's see.} \\
\\ 
\texttt{If I take 1 potato (2 points), then the other player can take the remaining 1 potato. That leaves 4 papers and 2 clocks. If I take some papers or clocks, I can get more points. For example, if I take 1 potato and 2 papers, that's 2 + 2 = 4 points. The other player could take 1 potato, 2 papers, and 2 clocks. But I don't know their values. If they value clocks more, they might prefer the clocks. Alternatively, if they value papers more, they might take the papers.} \\
\\ 
\texttt{Another option: take 1 potato and 1 clock. That's 2 + 1 = 3 points. Then the other player could take 1 potato, 1 clock, and the remaining 4 papers. But maybe that's not optimal for me. I need to maximize my own score while ensuring the other player can get their desired items.} \\
\\ 
\texttt{Wait, the problem says that if the proposals are complementary, meaning there are enough items to fulfill both, then we get our scores. Otherwise, zero. So my proposal must not require more of any item than available, and the other player's proposal must also fit within the remaining items.} \\
\\ 
\texttt{So, if I propose taking 2 potatoes, that's 2 potatoes. The other player's proposal must not require more than 0 potatoes. If they need any potatoes, that would conflict} \\
            }
        }
    }
    & & \\ \\

    \theutterance \stepcounter{utterance}  
    & & \multicolumn{4}{p{0.6\linewidth}}{
        \cellcolor[rgb]{0.9,0.9,0.9}{
            \makecell[{{p{\linewidth}}}]{
                \texttt{\tiny{[GM$\rangle$P2]}}
                \texttt{<think>} \\
\texttt{Okay, let's try to figure out the best way to approach this negotiation. My value function is paper:1, clock:1, potato:2. The other player's values are secret, so I need to make a proposal that's likely to be complementary. The key is to maximize my own score while ensuring that the other player can also get a set of items that doesn't overlap with mine.} \\
\\ 
\texttt{First, let's list all the items: 4 papers, 2 clocks, 2 potatoes. I need to decide how to divide these. Since I value potatoes the most, I should aim to get as many as possible. But there are only two potatoes. If I take both, that's 2*2=4 points. Then the remaining items are 4 papers and 2 clocks. The other player might want some of these. But if they also need potatoes, taking both could cause a conflict. So maybe I should take one potato and leave the other for them. Let's see.} \\
\\ 
\texttt{If I take 1 potato (2 points), then the other player can take the remaining 1 potato. That leaves 4 papers and 2 clocks. If I take some papers or clocks, I can get more points. For example, if I take 1 potato and 2 papers, that's 2 + 2 = 4 points. The other player could take 1 potato, 2 papers, and 2 clocks. But I don't know their values. If they value clocks more, they might prefer the clocks. Alternatively, if they value papers more, they might take the papers.} \\
\\ 
\texttt{Another option: take 1 potato and 1 clock. That's 2 + 1 = 3 points. Then the other player could take 1 potato, 1 clock, and the remaining 4 papers. But maybe that's not optimal for me. I need to maximize my own score while ensuring the other player can get their desired items.} \\
\\ 
\texttt{Wait, the problem says that if the proposals are complementary, meaning there are enough items to fulfill both, then we get our scores. Otherwise, zero. So my proposal must not require more of any item than available, and the other player's proposal must also fit within the remaining items.} \\
\\ 
\texttt{So, if I propose taking 2 potatoes, that's 2 potatoes. The other player's proposal must not require more than 0 potatoes. If they need any potatoes, that would conflict} \\
            }
        }
    }
    & \\ \\

    \theutterance \stepcounter{utterance}  
    & & & \multicolumn{4}{p{0.6\linewidth}}{
        \cellcolor[rgb]{0.9,0.9,0.9}{
            \makecell[{{p{\linewidth}}}]{
                \texttt{\tiny{[GM$\langle$P2]}}
                \texttt{<think>} \\
\texttt{Okay, let's try to figure out the best way to approach this negotiation. My value function is paper:1, clock:1, potato:2. The other player's values are secret, so I need to make a proposal that's likely to be complementary. The key is to maximize my own score while ensuring that the other player can also get a set of items that doesn't overlap with mine.} \\
\\ 
\texttt{First, let's list all the items: 4 papers, 2 clocks, 2 potatoes. I need to decide how to divide these. Since I value potatoes the most, I should aim to get as many as possible. But there are only two potatoes. If I take both, that's 2*2=4 points. Then the remaining items are 4 papers and 2 clocks. The other player might want some of these. But if they also need potatoes, taking both could cause a conflict. So maybe I should take one potato and leave the other for them. Let's see.} \\
\\ 
\texttt{If I take 1 potato (2 points), then the other player can take the remaining 1 potato. That leaves 4 papers and 2 clocks. If I take some papers or clocks, I can get more points. For example, if I take 1 potato and 2 papers, that's 2 + 2 = 4 points. The other player could take 1 potato, 2 papers, and 2 clocks. But I don't know their values. If they value clocks more, they might prefer the clocks. Alternatively, if they value papers more, they might take the papers.} \\
\\ 
\texttt{Another option: take 1 potato and 1 clock. That's 2 + 1 = 3 points. Then the other player could take 1 potato, 1 clock, and the remaining 4 papers. But maybe that's not optimal for me. I need to maximize my own score while ensuring the other player can get their desired items.} \\
\\ 
\texttt{Wait, the problem says that if the proposals are complementary, meaning there are enough items to fulfill both, then we get our scores. Otherwise, zero. So my proposal must not require more of any item than available, and the other player's proposal must also fit within the remaining items.} \\
\\ 
\texttt{So, if I propose taking 2 potatoes, that's 2 potatoes. The other player's proposal must not require more than 0 potatoes. If they need any potatoes, that would conflict} \\
            }
        }
    }
     \\ \\

    \theutterance \stepcounter{utterance}  
    & & & \multicolumn{2}{p{0.3\linewidth}}{
        \cellcolor[rgb]{0.9,0.9,0.9}{
            \makecell[{{p{\linewidth}}}]{
                \texttt{\tiny{[GM$|$GM]}}
                \texttt{continue} \\
            }
        }
    }
    & & \\ \\

    \theutterance \stepcounter{utterance}  
    & & & \multicolumn{2}{p{0.3\linewidth}}{
        \cellcolor[rgb]{0.9,0.9,0.9}{
            \makecell[{{p{\linewidth}}}]{
                \texttt{\tiny{[GM$|$GM]}}
                \texttt{<think>} \\
\texttt{Okay, let's try to figure out the best way to approach this negotiation. My value function is paper:1, clock:1, potato:2. The other player's values are secret, so I need to make a proposal that's likely to be complementary. The key is to maximize my own score while ensuring that the other player can also get a set of items that doesn't overlap with mine.} \\
\\ 
\texttt{First, let's list all the items: 4 papers, 2 clocks, 2 potatoes. I need to decide how to divide these. Since I value potatoes the most, I should aim to get as many as possible. But there are only two potatoes. If I take both, that's 2*2=4 points. Then the remaining items are 4 papers and 2 clocks. The other player might want some of these. But if they also need potatoes, taking both could cause a conflict. So maybe I should take one potato and leave the other for them. Let's see.} \\
\\ 
\texttt{If I take 1 potato (2 points), then the other player can take the remaining 1 potato. That leaves 4 papers and 2 clocks. If I take some papers or clocks, I can get more points. For example, if I take 1 potato and 2 papers, that's 2 + 2 = 4 points. The other player could take 1 potato, 2 papers, and 2 clocks. But I don't know their values. If they value clocks more, they might prefer the clocks. Alternatively, if they value papers more, they might take the papers.} \\
\\ 
\texttt{Another option: take 1 potato and 1 clock. That's 2 + 1 = 3 points. Then the other player could take 1 potato, 1 clock, and the remaining 4 papers. But maybe that's not optimal for me. I need to maximize my own score while ensuring the other player can get their desired items.} \\
\\ 
\texttt{Wait, the problem says that if the proposals are complementary, meaning there are enough items to fulfill both, then we get our scores. Otherwise, zero. So my proposal must not require more of any item than available, and the other player's proposal must also fit within the remaining items.} \\
\\ 
\texttt{So, if I propose taking 2 potatoes, that's 2 potatoes. The other player's proposal must not require more than 0 potatoes. If they need any potatoes, that would conflict} \\
            }
        }
    }
    & & \\ \\

    \theutterance \stepcounter{utterance}  
    & & \multicolumn{4}{p{0.6\linewidth}}{
        \cellcolor[rgb]{0.9,0.9,0.9}{
            \makecell[{{p{\linewidth}}}]{
                \texttt{\tiny{[P1$\langle$GM]}}
                \texttt{<think>} \\
\texttt{Okay, let's try to figure out the best way to approach this negotiation. My value function is paper:1, clock:1, potato:2. The other player's values are secret, so I need to make a proposal that's likely to be complementary. The key is to maximize my own score while ensuring that the other player can also get a set of items that doesn't overlap with mine.} \\
\\ 
\texttt{First, let's list all the items: 4 papers, 2 clocks, 2 potatoes. I need to decide how to divide these. Since I value potatoes the most, I should aim to get as many as possible. But there are only two potatoes. If I take both, that's 2*2=4 points. Then the remaining items are 4 papers and 2 clocks. The other player might want some of these. But if they also need potatoes, taking both could cause a conflict. So maybe I should take one potato and leave the other for them. Let's see.} \\
\\ 
\texttt{If I take 1 potato (2 points), then the other player can take the remaining 1 potato. That leaves 4 papers and 2 clocks. If I take some papers or clocks, I can get more points. For example, if I take 1 potato and 2 papers, that's 2 + 2 = 4 points. The other player could take 1 potato, 2 papers, and 2 clocks. But I don't know their values. If they value clocks more, they might prefer the clocks. Alternatively, if they value papers more, they might take the papers.} \\
\\ 
\texttt{Another option: take 1 potato and 1 clock. That's 2 + 1 = 3 points. Then the other player could take 1 potato, 1 clock, and the remaining 4 papers. But maybe that's not optimal for me. I need to maximize my own score while ensuring the other player can get their desired items.} \\
\\ 
\texttt{Wait, the problem says that if the proposals are complementary, meaning there are enough items to fulfill both, then we get our scores. Otherwise, zero. So my proposal must not require more of any item than available, and the other player's proposal must also fit within the remaining items.} \\
\\ 
\texttt{So, if I propose taking 2 potatoes, that's 2 potatoes. The other player's proposal must not require more than 0 potatoes. If they need any potatoes, that would conflict} \\
\\ 
\\ 
\texttt{The time is up. It is now your turn to submit a secret proposal.} \\
            }
        }
    }
    & \\ \\

    \theutterance \stepcounter{utterance}  
    & \multicolumn{4}{p{0.6\linewidth}}{
        \cellcolor[rgb]{0.9,0.9,0.9}{
            \makecell[{{p{\linewidth}}}]{
                \texttt{\tiny{[P1$\rangle$GM]}}
                \texttt{<think>} \\
\texttt{Okay, let's try to figure out the best way to approach this negotiation. My value function is paper:1, clock:1, potato:2. The other player's values are secret, so I need to make a proposal that's likely to be complementary. The key is to maximize my own score while ensuring that the other player can also get a set of items that doesn't overlap with mine.} \\
\\ 
\texttt{First, let's list all the items: 4 papers, 2 clocks, 2 potatoes. I need to decide how to divide these. Since I value potatoes the most, I should aim to get as many as possible. But there are only two potatoes. If I take both, that's 2*2=4 points. Then the remaining items are 4 papers and 2 clocks. The other player might want some of these. But if they also need potatoes, taking both could cause a conflict. So maybe I should take one potato and leave the other for them. Let's see.} \\
\\ 
\texttt{If I take 1 potato (2 points), then the other player can take the remaining 1 potato. That leaves 4 papers and 2 clocks. If I take some papers or clocks, I can get more points. For example, if I take 1 potato and 2 papers, that's 2 + 2 = 4 points. The other player could take 1 potato, 2 papers, and 2 clocks. But I don't know their values. If they value clocks more, they might prefer the clocks. Alternatively, if they value papers more, they might take the papers.} \\
\\ 
\texttt{Another option: take 1 potato and 1 clock. That's 2 + 1 = 3 points. Then the other player could take 1 potato, 1 clock, and the remaining 4 papers. But maybe that's not optimal for me. I need to maximize my own score while ensuring the other player can get their desired items.} \\
\\ 
\texttt{Wait, the problem says that if the proposals are complementary, meaning there are enough items to fulfill both, then we get our scores. Otherwise, zero. So my proposal must not require more of any item than available, and the other player's proposal must also fit within the remaining items.} \\
\\ 
\texttt{So, if I propose taking 2 potatoes, that's 2 potatoes. The other player's proposal must not require more than 0 potatoes. If they need any potatoes, that would conflict} \\
            }
        }
    }
    & & \\ \\

    \theutterance \stepcounter{utterance}  
    & & & \multicolumn{2}{p{0.3\linewidth}}{
        \cellcolor[rgb]{0.9,0.9,0.9}{
            \makecell[{{p{\linewidth}}}]{
                \texttt{\tiny{[GM$|$GM]}}
                \texttt{continue} \\
            }
        }
    }
    & & \\ \\

    \theutterance \stepcounter{utterance}  
    & & & \multicolumn{2}{p{0.3\linewidth}}{
        \cellcolor[rgb]{0.9,0.9,0.9}{
            \makecell[{{p{\linewidth}}}]{
                \texttt{\tiny{[GM$|$GM]}}
                \texttt{player was instructed to make a proposal} \\
            }
        }
    }
    & & \\ \\

\end{supertabular}
}

\end{document}
